\documentclass[11pt]{article}
% Text: the character-level half of the importer. Accents and ligatures
% become precomposed Unicode, the punctuation conventions (--- ``…'' ~) the
% characters they stand for, and the style commands become run flags,
% colours and sizes on the rich-document runs. What a Markdown view cannot
% show — a colour, a point size — is still carried in the model and lands in
% the ODT and DOCX exports.
\usepackage[T1]{fontenc}
\usepackage{xcolor}
\usepackage{ulem}
\usepackage{hyperref}
\definecolor{ultrablue}{RGB}{50,50,150}
\definecolor{warm}{HTML}{B4541E}
\colorlet{cool}{ultrablue}
\title{Text, Characters and Links}
\author{The UltraCanvas Team}
\date{September 2026}
\begin{document}
\maketitle

\begin{abstract}
Nothing in this document is a picture of text: every mark below is a
character the reader produced from its LaTeX spelling, which is what makes
the imported document searchable, selectable and translatable.
\end{abstract}

\section{Emphasis and style}
\textbf{Bold}, \textit{italic}, \underline{underlined}, \sout{struck out},
\textsc{Small Caps}, \textsf{sans serif} and \texttt{monospace}. The
declarations do the same for the rest of their group:
{\bfseries bold until the brace closes}, {\itshape italic until the brace
closes}.

\emph{Emphasis toggles rather than sets: inside an emphasised phrase,
\emph{a nested one goes back upright}, which is what \texttt{emph} means
and why it is not a synonym for \texttt{textit}.}

Sizes run from \tiny tiny\normalsize\ through \large large\normalsize\ to
\Huge huge\normalsize\ --- carried as a point size on the run.
Colours are named or mixed: \textcolor{ultrablue}{a defined RGB colour},
\textcolor{warm}{an HTML one}, \textcolor{cool}{a \texttt{colorlet} alias}
and \textcolor{red}{one of the dvips names}.

Subscripts and superscripts outside mathematics:
H\textsubscript{2}O, the 1\textsuperscript{st} of the month, and
E = mc\textsuperscript{2} written in text rather than as a formula.

\section{Characters}
\subsection{Accents}
Caf\'e, na\"{\i}ve, \^etre, Stra\ss e, se\~nor, \v{S}koda, Erd\H{o}s,
fa\c{c}ade, \k{a}\.z, \r{a}ngstr\"om, \u{g}\"unl\"uk --- each one a single
precomposed character where Unicode has it.

\subsection{Ligatures and symbols}
\ae sthetic, \oe uvre, \AA ngstr\"om, \O resund, \l od\'z, \DH, \TH, and the
text symbols \%, \$, \&, \_, \#, \{, \}.

\subsection{Punctuation}
A hyphen-like dash, an en dash 1914--1918, an em dash --- like this one ---
and an ellipsis\ldots{} ``Double quotes'' and `single quotes' curl the right
way round. A tie~binds two words, and \LaTeX{} and \TeX{} set their own
names.

\section{Links}
A labelled link to \href{https://ultra-os.org}{the project site}, a bare one
\url{https://code.claude.com/docs}, and \nolinkurl{example.com/not-a-link}
which is set as a URL without becoming one. A link keeps its target on the
run, so it is still a link after an export to Markdown, HTML, ODT or DOCX.

\section{Blocks of text}
\begin{itemize}
  \item An unordered item.
  \begin{enumerate}
    \item Ordered, nested one level.
    \item[(b)] An item with its own label.
  \end{enumerate}
  \item Back at the outer level.
\end{itemize}

\begin{quote}
A quotation is a block of its own, and it may hold \emph{any} run styling.
\end{quote}

\begin{center}
Centred text, as \texttt{center} and \texttt{\textbackslash centering} both
set it.
\end{center}

\begin{flushright}
Right-aligned text.
\end{flushright}

\hrule
\medskip
\noindent A horizontal rule above separates this closing paragraph, which is
set with \texttt{\textbackslash noindent} --- an instruction to a typesetter
that a rich document has no use for, and drops without a word.
\end{document}
